\lyseis{Α4.1}{Κανονικότητα, έννοια και αναπαραστάσεις}
\sol{1} α) Λ, Μ. Μονάδα επανάληψης: ΚΛΜ. β) $1,\ 1$. Μονάδα επανάληψης: $4,1,1$. γ) Β, Β. Μονάδα επανάληψης: ΑΒΒΓ.
\sol{2} α) $23,\ 28$ (κανόνας $+5$). β) $32,\ 25$ ($-7$). γ) $2,\ 5$ ($+3$). δ) $61,\ 72$ ($+11$).
\sol{3} α) $81,\ 243$ (πολλαπλασιάζω με $3$). β) $80,\ 160$ (με $2$). γ) $64,\ 32$ (κάθε όρος είναι το μισό του προηγούμενου, δηλαδή διαιρώ με $2$). δ) $32,\ -64$ (πολλαπλασιάζω με $-2$, γι' αυτό τα πρόσημα εναλλάσσονται).
\sol{4} α) Οι διαφορές είναι $1,2,3,4$, άρα επόμενες $5$ και $6$: $16$ και $22$. β) Οι διαφορές είναι $1,2,3,4$: επόμενα $17$ και $23$. γ) Οι διαφορές είναι $1,3,5,7$: επόμενες $9$ και $11$, άρα $25$ και $36$ (είναι τα τετράγωνα $0^2,1^2,2^2,\dots$).
\sol{5} α) Το $4$ο σχήμα είναι ένα «Γ» με $4$ τετράγωνα στην οριζόντια και $4$ στην κατακόρυφη πλευρά, με το γωνιακό κοινό: $7$ τετράγωνα (ακολουθία $1,3,5,7$). β) Το $4$ο σχήμα έχει $4$ γεμάτα τετράγωνα και ένα κενό: $5$ τετράγωνα (ακολουθία $2,3,4,5$).
\sol{6} α) Ο κανόνας είναι $+3$, άρα το τελευταίο πρέπει να είναι $13$. β) Ο κανόνας είναι $-10$, άρα $100,90,80,70,60$: το $75$ είναι λάθος. γ) Ο κανόνας είναι $\cdot2$, άρα το $49$ πρέπει να είναι $48$. δ) Το μοτίβο είναι ΚΜ ΚΜ ΚΜ\ldots Ο $5$ος όρος πρέπει να είναι Κ και όχι Μ.
\sol{7} α) Λάθος: υπάρχει πάντα κανόνας. β) Σωστό. γ) Λάθος: η μονάδα επανάληψης είναι η ΑΒ. δ) Σωστό ($+2$). ε) Σωστό ($+4$). στ) Σωστό. ζ) Λάθος: μπορούν να είναι και αρνητικοί αριθμοί ή σχήματα.
\sol{8} α) ΑΑΣΣΣ. β) Κάθε $5$ δευτερόλεπτα. γ) ΑΑΣΣΣ ΑΑΣΣΣ ΑΑΣΣΣ. δ) Τα δευτερόλεπτα $16$ και $17$ είναι τα δύο πρώτα της $4$ης μονάδας, δηλαδή ΑΑ: τα λαμπάκια είναι αναμμένα.
\sol{9} α) Οι χρόνοι είναι $4,5,6,7,8,9$: το $6$ο σκαλί παίρνει $9$ δευτερόλεπτα. β) Το $4$ είναι στη θέση $1$ και το $15$ είναι $11$ μονάδες μετά, άρα στη θέση $12$: το $12$ο σκαλί.
\sol{10} Οι κωδικοί διπλασιάζονται: $3,6,12,24,48,96,192$, άρα ο $7$ος είναι ο $192$. Συνεχίζουμε: $384,\ 768$. Ο πρώτος που ξεπερνά το $500$ είναι ο $768$ (ο $9$ος).
\sol{11} α) Τα πλακάκια είναι $1,3,6,10,15$: το $4$ο έχει $10$ και το $5$ο έχει $15$ (προσθέτω κάθε φορά $2,3,4,5$). β) Η περίμετρος αυξάνεται κατά $4$ κάθε φορά: $16$ για το $4$ο και $20$ για το $5$ο. Κανόνας: $+4$.
\sol{12} α) Κάθε όρος είναι κατά $3$ μεγαλύτερος από τον προηγούμενο ($+3$). β) Τα σημεία είναι τα $(1,2)$, $(2,5)$, $(3,8)$, $(4,11)$: ανεβαίνουν κατά $3$ όταν το $\nu$ αυξάνεται κατά $1$. γ) $5$ος όρος: $14$, $6$ος όρος: $17$.
\sol{13} Και οι δύο μπορεί να έχουν δίκιο: αν ο κανόνας είναι «διπλασιάζω», έχουμε $2,4,8,16$. Αν ο κανόνας είναι «οι διαφορές αυξάνονται κατά $2$» ($2,\,4,\,6,\dots$), έχουμε $2,4,8,14$. Τρεις όροι δεν αρκούν για να ξεχωρίσουμε τον κανόνα.
\sol{14} Μετά από $7$ ημέρες ξανά Δευτέρα. Η $1$η, $8$η, $15$η, \ldots είναι Δευτέρες. Η $99$η ($=1+14\cdot7$) είναι Δευτέρα, άρα η $100$ή είναι Τρίτη.
\sol{15} Μονάδα επανάληψης: ΑΒΑ ($3$ γράμματα). Στα $12$ γράμματα έχουμε $4$ μονάδες και κάθε μονάδα έχει $2$ Α, άρα $8$ Α.
\sol{16} Οι διαφορές είναι $1,2,3,\dots,9$. Οι όροι είναι $1,2,4,7,11,16,22,29,37,46$. Ο $10$ος όρος είναι $46$.
\sol{17} Ελεύθερη απάντηση. Π.χ. οι εποχές (κανόνας: ξαναρχίζουν κάθε $4$), η σειρά των φαναριών (πράσινο, πορτοκαλί, κόκκινο), οι ημέρες της εβδομάδας.
\sol{18} Ελεύθερη κατασκευή. Ο κανόνας πρέπει να ισχύει σε όλους τους όρους και να επιτρέπει να προβλέψεις τους επόμενους.
